% concatenated from neurips_2026.tex, sections/abstract.tex, sections/intro.tex, sections/problem.tex, sections/bg.tex, sections/grad-approx.tex, sections/exp.tex, sections/fut.tex, sections/ack.tex, sections/appendix.tex

\documentclass{article}

% if you need to pass options to natbib, use, e.g.:
 \PassOptionsToPackage{numbers, compress}{natbib}
% before loading neurips_2026

% The authors should use one of these tracks.
% Before accepting by the NeurIPS conference, select one of the options below.
% 0. "default" for submission
%\usepackage{neurips_2026}  %For conference
\usepackage{amsmath}   % equation environments, \operatorname, \text, \underbrace
\usepackage{amsthm}    % theorem, lemma, proposition, proof environments
\usepackage{amssymb}   % \mathbb, \odot
\usepackage{empheq}    % for \boxed inside equation environments
\newtheorem{theorem}{Theorem}
\newtheorem{lemma}[theorem]{Lemma}
\newtheorem{proposition}[theorem]{Proposition}
\usepackage{subcaption}
% the "default" option is equal to the "main" option, which is used for the Main Track with double-blind reviewing.
% 1. "main" option is used for the Main Track
%  \usepackage[main]{neurips_2026}
% 2. "position" option is used for the Position Paper Track
%  \usepackage[position]{neurips_2026}
% 3. "eandd" option is used for the Evaluations & Datasets Track
 % \usepackage[eandd]{neurips_2026}
 % if you want to opt-in for a de-anomymized submission (not recommended, but OK):
 %\usepackage[eandd, nonanonymous]{neurips_2026}
% 4. "creativeai" option is used for the Creative AI Track
%  \usepackage[creativeai]{neurips_2026}
% 5. "sglblindworkshop" option is used for the Workshop with single-blind reviewing
 % \usepackage[sglblindworkshop]{neurips_2026}
% 6. "dblblindworkshop" option is used for the Workshop with double-blind reviewing
%  \usepackage[dblblindworkshop]{neurips_2026}

% After being accepted, the authors should add "final" behind the track to compile a camera-ready version.
% 1. Main Track
 % \usepackage[main, final]{neurips_2026}
% 2. Position Paper Track
%  \usepackage[position, final]{neurips_2026}
% 3. Evaluations & Datasets Track
 % \usepackage[eandd, final]{neurips_2026}
% 4. Creative AI Track
%  \usepackage[creativeai, final]{neurips_2026}
% 5. Workshop with single-blind reviewing
%  \usepackage[sglblindworkshop, final]{neurips_2026}
% 6. Workshop with double-blind reviewing
%  \usepackage[dblblindworkshop, final]{neurips_2026}
% Note. For the workshop paper template, both \title{} and \workshoptitle{} are required, with the former indicating the paper title shown in the title and the latter indicating the workshop title displayed in the footnote.
% For workshops (5., 6.), the authors should add the name of the workshop, "\workshoptitle" command is used to set the workshop title.
% \workshoptitle{WORKSHOP TITLE}

% "preprint" option is used for arXiv or other preprint submissions
  \usepackage[preprint]{neurips_2026}

% to avoid loading the natbib package, add option nonatbib:
%    \usepackage[nonatbib]{neurips_2026}

\usepackage[utf8]{inputenc} % allow utf-8 input
\usepackage[T1]{fontenc}    % use 8-bit T1 fonts
\usepackage{hyperref}       % hyperlinks
\usepackage{url}            % simple URL typesetting
\usepackage{booktabs}       % professional-quality tables
\usepackage{amsfonts}       % blackboard math symbols
\usepackage{nicefrac}       % compact symbols for 1/2, etc.
\usepackage{microtype}      % microtypography
\usepackage{xcolor}         % colors
\usepackage{algorithm}
\usepackage{algpseudocode}
\usepackage{float}

% Note. For the workshop paper template, both \title{} and \workshoptitle{} are required, with the former indicating the paper title shown in the title and the latter indicating the workshop title displayed in the footnote. 
\title{Fast and Stable Gradient Approximation for Bilinear Forms of Hermitian Matrix Functions}


% The \author macro works with any number of authors. There are two commands
% used to separate the names and addresses of multiple authors: \And and \AND.
%
% Using \And between authors leaves it to LaTeX to determine where to break the
% lines. Using \AND forces a line break at that point. So, if LaTeX puts 3 of 4
% authors names on the first line, and the last on the second line, try using
% \AND instead of \And before the third author name.


\author{%
  Navjot Singh \\
  Lawrence Berkeley National Laboratory\\
  Berkeley, CA 94720 \\
  \texttt{nsingh2@lbl.gov} \\
  % examples of more authors
  \And
  Kipton Barros \\
  % Theoretical Division and CNLS \\
  Los Alamos National Laboratory \\
  Los Alamos, NM 87545 \\
  \texttt{kbarros@lanl.gov} \\
  \And
  Xiaoye Sherry Li\\
  Lawrence Berkeley National Laboratory\\
  Berkeley, CA 94720 \\
  \texttt{xsli@lbl.gov} 
}


\begin{document}


\maketitle

% Abstract
\begin{abstract}
Objectives involving bilinear forms \(u^\top f(A(\theta))v\) for Hermitian \(A\) arise widely in scientific computing and probabilistic machine learning. For large matrices, Lanczos efficiently approximates these quantities, but differentiating them with respect to \(\theta\) is challenging. Existing approaches either backpropagate through the Lanczos recurrence, requiring reorthogonalization for stability, or apply Arnoldi to an augmented block matrix of twice the original size. Both introduce extra computation and orthogonalization costs that can limit performance on modern hardware.
We propose a forward-only gradient approximation that reuses the Lanczos pass and adds 
very minimal overhead in most cases.
% only matrix--matrix products and computations with the small tridiagonal matrix.
We prove that its error is proportional to the Lanczos residual norm, the same quantity controlling the forward approximation. 
% Experiments show that the method remains stable without reorthogonalization where adjoint-based alternatives become unstable, and is faster than existing state-of-the-art approaches.
Whereas a traditional adjoint-based calculation would be unstable without reorthogonalization, the new method appears unconditionally stable in our tests. It is also faster than existing state-of-the-art approaches.
\end{abstract}


\section{Introduction}
Large-scale learning, inverse problems, and scientific computing involve objectives built from \emph{functions of 
%parameter-dependent
parameterized matrices}, rather than from individual matrix entries. Examples include log-determinants in Gaussian processes and marginal-likelihood optimization, matrix exponentials in dynamical systems and network analysis, and time-evolution operators in quantum models \citep{rasmussen2006gaussian,dong2017scalable,immer2021scalable,saad1992exponential,hochbruck1997krylov,benzi2020,banchi2023learning}. In these applications, the matrix \(A(\theta)\) is often too large to form, diagonalize, or differentiate as a dense object, but matrix--vector products \(v\mapsto A(\theta)v\) and their differentials are available through sparse linear algebra, structured operators, or automatic differentiation. Krylov methods, and in particular the Lanczos method for Hermitian matrices, are therefore a natural computational primitive for matrix-function actions and quadratic forms \citep{gallopoulos1992parabolic,beckermann2009faber,ubaru2017fast,chen2023krylov,davies2021krylov}: an \(m\)-step Lanczos iteration approximates quadratic forms $u^\top f(A(\theta))u$
by replacing the large matrix function with a small tridiagonal matrix function. Differentiating this approximation, however, is not straightforward. Exact Fr\'echet-derivative formulas are well understood \citep{najfeld1995,almohy2009}, and recent Krylov methods approximate Fr\'echet-derivative actions without forming the full derivative \citep{kressner2019,kandolf2023,crouzeix2020,kressner2026}. Alternatively, one may differentiate through the Lanczos iterations using reverse-mode automatic differentiation or custom adjoint recurrences. Although these approaches are principled, 
%but they either introduce an additional derivative Krylov problem or require a backward pass through the Lanczos recurrence, whose finite-precision behavior can be sensitive to loss of orthogonality.
they do not inherit the nice numerical stability properties of the forward Lanczos quadrature calculation. In finite numerical precision, such instability of the gradient calculation may force the use of an expensive reorthogonalization procedure.

We propose a forward-only approximate gradient for
\[
    \phi(\theta;u)=u^\top f(A(\theta))u,
\]
where \(A(\theta)\) is Hermitian and the forward value is computed by Lanczos. The central idea is to differentiate only the small projected Lanczos problem, and then lift this sensitivity back to the parameters through vector--Jacobian products of the original matrix--vector product map \((\theta,v)\mapsto A(\theta)v\). Our contributions are:
\begin{itemize}
    \item We derive an approximate gradient for Lanczos quadratic-form approximations that requires only the forward Lanczos calculation, some computation involving the small tridiagonal matrix, and the usual differentiation of matrix-vector products.
    \item We prove that Lanczos basis-variation terms can be ignored up to an error governed by the final Lanczos residual coefficient \(\beta_m\), which therefore vanishes when the Krylov subspace is invariant.
    \item We show empirically that the estimator inherits the numerical stability of the forward Lanczos approximation, remaining reliable without full reorthogonalization in settings where naive backpropagation or adjoint-based methods break down.
    \item We demonstrate the approach on log-determinant, network-sensitivity, and Hamiltonian-learning experiments, where it improves upon the accuracy or matches it while being more efficient.
\end{itemize}

The rest of the paper is organized as follows. Section~\ref{sec:problem} formulates the problem setup, establishing the relationship between other objectives and the quadratic form, as well as the model assumptions. Section~\ref{sec:background} provides brief background on the methods used to tackle the problem. Section~\ref{sec:grad_approx} contains our main contributions and derives the gradient approximation. Section~\ref{sec:experiments} evaluates our method against various benchmarks.



\section{Problem setup and model assumption}
\label{sec:problem}
\subsection{Problem setup}

Let $\theta \in \Theta$ denote the model parameters, and let $
A(\theta) \in \mathbb{R}^{n \times n}$ be a symmetric matrix depending smoothly on $\theta$. All the proofs and algorithms in this manuscript generalize readily to $A(\theta) \in \mathbb{C}^{n \times n}$ Hermitian matrices, but we restrict to real
matrices and vectors for simplicity. We assume that the spectrum of $A(\theta)$ is contained in an interval
$\mathcal I \subset \mathbb{R}$ for all $\theta$ of interest, and that
\begin{equation*}
f : \mathcal I \to \mathbb{R}
\end{equation*}
is a scalar function for which the matrix function $f(A(\theta))$ is well defined.

 % Comment: Add we consider real, everything acn be made complex
 


Our primary object of interest is the quadratic form
\begin{equation}
\phi(\theta;u) := u^\top f(A(\theta))u,
\qquad u \in \mathbb{R}^n.
\label{eq:problem_quadratic_form}
\end{equation}
We focus on \eqref{eq:problem_quadratic_form} because it is a basic building block for several
quantities of interest involving matrix functions. In particular, trace objectives can be written as
expectations of quadratic forms,
\begin{equation}
\operatorname{tr}(f(A(\theta)))
= \mathbb E_z\!\left[z^\top f(A(\theta)) z\right],
\qquad \mathbb E[zz^\top]=I,
\label{eq:problem_trace_identity}
\end{equation}
for stochastic $z$~\citep{hutchinson1989stochastic}. Bilinear forms can be recovered from quadratic forms by polarization,
\begin{equation}
u^\top f(A(\theta))v
=
\frac14\Big((u+v)^\top f(A(\theta))(u+v)
-
(u-v)^\top f(A(\theta))(u-v)\Big).
\label{eq:problem_polarization}
\end{equation}
Thus, once we can compute gradients of \eqref{eq:problem_quadratic_form}, the same machinery also
extends to stochastic trace estimation and more general objectives involving bilinear forms. Our goal is therefore to compute the gradient $\nabla_\theta \phi(\theta;u)$  with respect to $\theta$.


\subsection{Matrix-free access model}

We are interested in the large-scale regime in which $A(\theta)$ is too expensive to
form densely or factorize explicitly. Instead, we assume matrix-free access through
differentiable matrix--vector products:
\begin{equation}
    \mathcal{A}(\theta,u) = A(\theta)u .
    \label{eq:problem_matvec_access}
\end{equation}
This setting covers sparse and structured operators arising from discretized PDEs and
other scientific-computing problems \citep{saad2003iterative,elman2014finite},
precision and covariance operators in spatial statistics and Gaussian-process models
\citep{rue2005gaussian,rasmussen2006gaussian}, generalized Gauss--Newton and other
curvature matrices in large-scale optimization \citep{schraudolph2002fast,martens2010deep},
and more generally any linear operator for which matrix--vector products are inexpensive
even when $n$ is large.

For gradient computation, we additionally assume that the map
$\mathcal{A}(\theta,u)$ in \eqref{eq:problem_matvec_access} is differentiable with
respect to $\theta$. Thus, for fixed $u$, we can evaluate Jacobian--vector or
vector--Jacobian products associated with $\theta\mapsto A(\theta)u$. Equivalently, we
assume access to differential matrix--vector products without materializing
$\partial A(\theta)/\partial\theta$.

\section{Background}
\label{sec:background}
\subsection{Lanczos approximation of matrix functions}
\label{subsec:bg_lanczos}

Lanczos iteration is a standard matrix-free tool for approximating
matrix-function-vector products \(f(A)u\) when \(A\) is large and symmetric
\citep{golub2013matrix,higham2008functions}. Starting from \(v_1=u/\|u\|\),
\(m\) Lanczos steps construct an orthonormal basis \(V=[v_1,\ldots,v_m]\) for
the Krylov subspace
\[
    \mathcal{K}_m(A,u)=\operatorname{span}\{u,Au,\ldots,A^{m-1}u\},
\]
together with a symmetric tridiagonal matrix \(T\) satisfying
\begin{equation}
    A V = V T + \beta_m v_{m+1} e_m^\top .
    \label{eq:bg_lanczos_relation}
\end{equation}
The resulting Krylov approximation is
\begin{equation}
    f(A)u
    \approx
    \|u\| V f(T)e_1 ,
    \label{eq:bg_lanczos_mvf}
\end{equation}
so the nonlinear function is evaluated only on the small projected matrix \(T\). The
same projection gives the quadratic-form approximation
\begin{equation}
    u^\top f(A)u
    \approx
    \|u\|^2 e_1^\top f(T)e_1 .
    \label{eq:bg_lanczos_quadratic}
\end{equation}
In exact arithmetic, the Lanczos vectors are mutually orthogonal and \(T\)
contains no repeated Ritz values. In finite precision, however, the three-term
recurrence may lose orthogonality as Ritz values converge, which can lead to
``ghost'' copies of already converged eigenvalues in the $T$ matrix
\citep{paige1971computation,paige1972computational,paige1976error,parlett1998symmetric}.
This loss of orthogonality does not necessarily make Lanczos-based quadrature
unstable: the computed recurrence can often still be interpreted as a valid
Gauss quadrature rule for a nearby measure, and the associated quadratic-form
estimates are typically robust \citep{golub2009matrices,meurant2006lanczos}.
Nevertheless, ghost Ritz values can slow or distort convergence. Full
reorthogonalization restores stability but significantly increases the cost of each Lanczos
step as a QR decomposition with cost $O(nk^2)$ is needed at the $k$th iteration; cheaper alternatives include selective or partial reorthogonalization,
which maintain orthogonality only to the extent needed to suppress spurious
copies of converged Ritz values \citep{parlett1979lanczos,simons1984analysis}.

When \(A=A(\theta)\), one direct way to compute gradients is to differentiate the
finite Lanczos approximation itself. In automatic differentiation frameworks, this
amounts to backpropagating through the Krylov iteration. Kramer et
al.~\citep{kramer2024gradients} show that this can be inefficient for Lanczos and
Arnoldi iterations, especially when stable implementations require reorthogonalization.


\subsection{Gradients of matrix functions}
\label{subsec:bg_matrix_function_gradients}

Gradients of objectives involving \(f(A(\theta))\) are naturally expressed using
Fr\'echet derivatives~\citep{higham2008functions}. For a matrix function \(f\), the Fr\'echet derivative at \(A\) is
the linear map \(L_f(A,\cdot)\) defined by 
\begin{equation}
    f(A+\varepsilon E)
    =
    f(A)+\varepsilon L_f(A,E)+o(\varepsilon),
    \qquad \varepsilon\to 0 .
    \label{eq:bg_frechet_def}
\end{equation}
For symmetric \(A=Q\Lambda Q^\top\), it admits the divided-difference representation
\citep{higham2008functions}
\begin{equation}
    L_f(A,E)
    =
    Q\bigl(F(\Lambda)\circ (Q^\top E Q)\bigr)Q^\top ,
    \label{eq:bg_hadamard_frechet}
\end{equation}
where $\circ$ is the elementwise Hadamard product, and
\begin{align}
    F(\lambda_i,\lambda_j)
    =
    \begin{cases}
        \dfrac{f(\lambda_i)-f(\lambda_j)}{\lambda_i-\lambda_j}, & i\neq j,\\[1.2ex]
        f'(\lambda_i), & i=j.
    \end{cases}
\label{eq:bg_divided_difference}
\end{align}
This formula is useful analytically, but it requires spectral information about the full
matrix and dense \(n\times n\) quantities, making it impractical as a large-scale
algorithm.

A common matrix-free alternative is based on the block triangular identity
\citep{higham2008functions,najfeld1995derivatives}
\begin{equation}
    f\!\left(
    \begin{bmatrix}
        A & E\\
        0 & A
    \end{bmatrix}
    \right)
    \begin{bmatrix}
        0\\ u
    \end{bmatrix}
    =
    \begin{bmatrix}
        L_f(A,E)u\\
        f(A)u
    \end{bmatrix}.
    \label{eq:bg_block_action}
\end{equation}
For \(E=\partial A(\theta)/\partial\theta\), this computes the derivative action
\(L_f(A,E)u\) and the primal action \(f(A)u\) simultaneously. However, even if \(A\) is
symmetric, the augmented matrix is generally nonsymmetric and may have less favorable
spectral properties, so the Lanczos structure is lost and one typically uses Arnoldi.
Kressner and Oehme address this issue by modifying Arnoldi to better preserve the block
triangular structure~\citep{kressner2026}.

Another approach is to define a custom reverse-mode rule for the Lanczos iteration.
Kramer et al.~\citep{kramer2024gradients} formulate Lanczos and Arnoldi iterations as
algebraic constraints and apply the adjoint method to those constraints. For a scalar
loss $\rho$ depending on the Lanczos outputs, including the Lanczos vectors and the
tridiagonal coefficients, their adjoint system takes as input the corresponding output
sensitivities and solves a backward recurrence for adjoint variables. The resulting matrix
gradient has the form
\begin{equation}
    \nabla_A \rho
    =
    \sum_{j=1}^m \lambda_j v_j^\top ,
    \label{eq:bg_lanczos_adjoint_matrix_gradient}
\end{equation}
where the adjoint vectors $\lambda_j$ are obtained from the backward recurrence. If
$A=A(\theta)$ is accessed through differentiable matrix-vector products, parameter
gradients are then computed by contracting this matrix gradient with
$\partial_\theta A(\theta)$, or equivalently through vector-Jacobian products of the
matrix-vector product map.

For the Lanczos approximation
of the quadratic objective, the required output sensitivities include the
derivatives of the small projected objective with respect to the diagonal and
off-diagonal entries of $T$. These sensitivities are obtained by applying the
Fr\'echet derivative formula to the scalar map
$T\mapsto \|u\|^2 e_1^\top f(T)e_1$, and they initialize the backward Lanczos
recurrence.

Adjoint methods avoid generic backpropagation through the implemented loop, but they
still differentiate the Lanczos process by running a backward adjoint recurrence. This
motivates the question we address next: whether the projected sensitivity of the small
Lanczos objective can be used without differentiating through the Lanczos recurrence
itself.

\section{Gradient approximation of quadratic forms of matrix functions}
\label{sec:grad_approx}
We derive an efficient approximate gradient of the Lanczos estimate
\begin{equation}
    \widehat{\phi}(\theta;u)
    =
    \|u\|^2 e_1^\top f(T)e_1
    \label{eq:method_lanczos_quad}
\end{equation}
with respect to $\theta$, without backpropagating through the Lanczos recurrence.
Here $T$ is the tridiagonal matrix produced by running $m$ steps of Lanczos on
$A(\theta)$ with initial vector $v_1=u/\|u\|$. For clarity, we treat $\theta$ as a
scalar parameter; the extension to multiple parameters follows trivially. Throughout, we write $A=A(\theta)$, and all differentials are
with respect to $\theta$.

\subsection{Sensitivity with respect to the projected matrix}
\label{subsec:method_T_sensitivity}
\label{sec:sensitivity}

The first ingredient is the derivative of the scalar projected quantity with respect to
its small matrix argument. This is a standard consequence of the Fr\'echet derivative of
a matrix function.

\begin{theorem}[Projected sensitivity]
\label{thm:projected_sensitivity}
Let $T\in\mathbb{R}^{m\times m}$ be symmetric with eigendecomposition
$T=Q\Lambda Q^\top$, let $c=Q^\top e_1$, and let $F=F(\Lambda)$ be the
divided-difference matrix from \eqref{eq:bg_divided_difference}. For
\[
    \widehat{\phi}(T)=\|u\|^2 e_1^\top f(T)e_1,
\]
the differential with respect to $T$ is
\begin{equation}
    d\widehat{\phi}
    =
    \operatorname{tr}(G^\top dT),
    \label{eq:method_dphi_G}
\end{equation}
where
\begin{equation}
    G
    =
    \|u\|^2 Q\bigl((cc^\top)\circ F\bigr)Q^\top .
    \label{eq:method_G}
\end{equation}
\end{theorem}

Theorem~\ref{thm:projected_sensitivity} states that $G\in\mathbb{R}^{m\times m}$ is the
sensitivity of the scalar Lanczos estimate with respect to perturbations of the tridiagonal
matrix $T$ and is symmetric. The full derivation is given in Appendix~\ref{app:sensitivity}.

\subsection{Gradient approximation}
\label{subsec:method_gradient_approx}
\label{sec:gm}

Since $T=V^\top A V$, its full differential is
\begin{equation}
    dT
    =
    \underbrace{V^\top dA\,V}_{\text{direct}}
    +
    \underbrace{dV^\top A V + V^\top A\,dV}_{\text{basis variation}} .
    \label{eq:method_dT_full}
\end{equation}
The direct term captures the explicit dependence of $T$ on $A$, while the remaining
terms capture the dependence of the Lanczos basis $V$ on $\theta$. We approximate
\begin{equation}
    dT \approx V^\top dA\,V,
    \label{eq:method_dT_approx}
\end{equation}
dropping the basis variation terms. Substituting \eqref{eq:method_dT_approx} into \eqref{eq:method_dphi_G} yields
\begin{equation}
    d\widehat{\phi}
    \approx
    \operatorname{tr}\!\left(G V^\top dA\,V\right)
    \label{eq:method_grad_approx}
\end{equation}

The following theorem shows that the terms omitted in \eqref{eq:method_dT_approx} are
controlled by the Lanczos residual coefficient.

\begin{theorem}[Error from ignoring basis variation]
\label{thm:method_error}
Let $V,T,\beta_m,v_{m+1}$ satisfy the Lanczos relation
\eqref{eq:bg_lanczos_relation}, and let $G$ be defined by \eqref{eq:method_G}. Write
the differential of the Lanczos basis as
\[
    dV = VS + V^\perp N,
\]
where $S$ is skew-symmetric, $V^\perp$ spans the orthogonal complement of
$\operatorname{span}(V)$, and the first column of $V^\perp$ is $v_{m+1}$. Define
\[
    \eta = N^\top e_1^{(n-m)} .
\]
Then the exact differential of the Lanczos estimate satisfies
\begin{equation}
    d\widehat{\phi}
    =
    \operatorname{tr}\!\left(VGV^\top dA\right)
    +
    2\beta_m e_m^\top G\eta .
    \label{eq:method_error_formula}
\end{equation}
Consequently, the error in \eqref{eq:method_grad_approx} is proportional to the
Lanczos residual coefficient $\beta_m$ and vanishes when the Krylov subspace is
$A$-invariant.
\end{theorem}

\begin{proof}[Proof sketch]
Differentiating $T=V^\top A V$ gives \eqref{eq:method_dT_full}. Since $V^\top V=I$,
the basis differential decomposes as $dV=VS+V^\perp N$ with $S^\top=-S$. Using the
Lanczos relation
\[
    AV = VT+\beta_m v_{m+1}e_m^\top,
\]
one obtains
\[
    dT
    =
    V^\top dA\,V
    +
    [T,S]
    +
    \beta_m(\eta e_m^\top+e_m\eta^\top),
\]
where $[T,S]=TS-ST$. Substituting this expression into
$d\widehat{\phi}=\operatorname{tr}(G^\top dT)$ gives a direct term, a commutator term,
and a boundary term. The commutator term vanishes because the initial Lanczos vector
$v_1=u/\|u\|$ is fixed, which implies $Se_1=0$. The remaining boundary term equals
$2\beta_m e_m^\top G\eta$ since $G$ is symmetric. Full details are given in
Appendix~\ref{app:error}.
\end{proof}


\subsection{Practical implementation}
\label{subsec:method_implementation}

The approximation in \eqref{eq:method_grad_approx} can be evaluated without forming
the dense matrix $VGV^\top$. We present a practical algorithm below to compute the gradients through Lanczos iteration when $dA$ is not given as an explicit matrix.

\paragraph{Step 1: Forward Lanczos pass.}
Run $m$ steps of Lanczos with $A(\theta)$ and initial vector $v_1=u/\|u\|$ to obtain
\[
    V=[v_1,\ldots,v_m]\in\mathbb{R}^{n\times m},
    \qquad
    T\in\mathbb{R}^{m\times m}.
\]

\paragraph{Step 2: Projected sensitivity.}
Compute $T=Q\Lambda Q^\top$, set $c=Q^\top e_1$, form the divided-difference matrix
$F$, and construct
\[
    G
    =
    \|u\|^2 Q\bigl((cc^\top)\circ F\bigr)Q^\top .
\]

\paragraph{Step 3: parameter contractions.}
Let
\[
    W = VG = [w_1,\ldots,w_m].
\]
Then \eqref{eq:method_grad_approx} can be written as
\begin{equation}
    \frac{d\widehat{\phi}}{d\theta}
    \approx
    \sum_{j=1}^{m}
    w_j^\top
    \frac{\partial A(\theta)}{\partial\theta}
    v_j .
    \label{eq:method_grad_vjp}
\end{equation}
For fixed $v_j$, $w_j$ can be applied to the Jacobian
$D_\theta\mathcal{A}(\theta,v_j)$ 
\emph{vector-Jacobian product} (VJP)
$ w_j^\top [D_\theta\mathcal{A}(\theta,v_j)],
$
and its value is exactly the $j$th scalar contraction in
\eqref{eq:method_grad_vjp}. Thus the backward computation requires only $m$ such VJPs
through the same matrix-vector product primitive used in the forward Lanczos pass, and
does not require unrolling or differentiating through the Lanczos recurrence. 

We summarize the above details in Algorithm~\ref{alg:method_gradient}.


\begin{algorithm}[H]
\caption{Approximate gradient of $\phi(\theta;u)=u^\top f(A(\theta))u$}
\label{alg:method_gradient}
\begin{algorithmic}[1]
\Require Parameter $\theta$, vector $u$, Lanczos steps $m$, function $f$ and derivative $f'$
\Ensure Approximate gradient $\nabla_{\theta} \phi(\theta ; u)$
\State Run $m$-step Lanczos on $A(\theta)$ with $v_1=u/\|u\|$ to obtain $V$ and $T$
\State Compute $T=Q\Lambda Q^\top$ and set $c\leftarrow Q^\top e_1$
\State Form the divided-difference matrix $F$ using $f$ and $f'$ as in~\eqref{eq:bg_divided_difference} and form
 \[G \leftarrow \|u\|^2 Q\bigl((cc^\top)\circ F\bigr)Q^\top\]
 \State Compute the product \[W \leftarrow VG\]
\State \Return $\displaystyle
    \sum_{j=1}^{m}
    w_j^\top[D_\theta\mathcal{A}(\theta,v_j)]$
\end{algorithmic}
\end{algorithm}

\paragraph{Complexity and memory.}
Assuming $m$ Lanczos iterations, the dominant large-scale costs are the $m$ matrix-vector products in the forward
Lanczos pass and the $m$ VJPs in \eqref{eq:method_grad_vjp}. The remaining operations
involve only the stored basis and the small projected matrix: The symmetric tridiagonal
eigendecomposition of $T$ costs $O(m^2)$. The dense matrix-matrix multiply $W = VG$ costs $O(nm^2)$ with very low prefactor.
Forming $G$ directly costs $O(m^3)$.
An alternative is to construct $G V$ directly at $O(nm^2)$ cost using optimized dense linear algebra kernels.
Similarly, the adjoint products in \eqref{eq:method_grad_vjp} are independent across
Lanczos steps and can be parallelized or batched when the matrix-vector product
primitive supports it. 



\section{Experiments}
\label{sec:experiments}
Our experiments test whether the proposed method gives accurate and efficient
gradients for objectives that reduce to quadratic forms of matrix
functions. We consider three settings: $u^\top \log(K)u$, where we compare with
adjoint-based Lanczos differentiation used in~\cite{kramer2024gradients}; graph-sensitivity objectives involving
$u^\top \exp(A)v$, where we compare with a one-pass block Arnoldi
Fr\'echet-derivative approximation used in~\cite{kressner2026}; and an end-to-end Hamiltonian-learning task,
where we compare with a dense Fr\'echet-derivative reference. All experiments were run on a MacBook Pro with an Apple M4 Pro chip
(12-core CPU: 8 performance cores and 4 efficiency cores), 24 GB unified memory,
running macOS 26.2 on arm64.

\paragraph{Implementation details.}
To ensure fairness, for the logarithm experiments we implement our forward-only gradient
estimator within the same \texttt{matfree}-based codebase used by the Lanczos-adjoint
baseline of \citet{kramer2024gradients}. Thus, our method and the adjoint baseline share
the same Lanczos forward pass, matrix-vector-product interface, objectives, and reference
routines; only the gradient computation differs. For the block-Arnoldi comparisons, we
use the public \texttt{fAb}-Frechet repository accompanying \citet{kressner2026}, reusing
its graph-loading code, network-sensitivity objectives, and reference computations.
Sparse matrix-vector products are evaluated with the NumPy/SciPy sparse routines used
by that implementation.



\subsection{Matrix-logarithm and their derivatives}
\label{subsec:experiments_logdet}

Log-determinants of symmetric positive definite matrices are a standard bottleneck in
Gaussian-process marginal likelihoods and related models involving covariance, kernel,
precision, Hessian, or Gauss--Newton matrices
\citep{rasmussen2006gaussian,gardner2018gpytorch,ubaru2017fast,dong2017scalable,immer2021scalable}.
For an SPD matrix $K(\theta)$,
\[
    \log\det K(\theta)=\operatorname{tr}(\log K(\theta)),
\]
and stochastic Lanczos quadrature estimates this trace by averaging quadratic forms
$u^\top \log(K(\theta))u$ over random probe vectors $u$. In this experiment, we isolate
a single fixed Rademacher probe. This removes stochastic trace-estimation variance and
tests the core question: whether the gradient of the Lanczos approximation to
$u^\top\log(K(\theta))u$ is accurate.

We construct $K(\theta)$ from a synthetic Gaussian-process regression problem with
inputs $x_i\in\mathbb{R}^2$ and an RBF kernel with diagonal noise,
\[
K_{ij}(\theta)
=
\sigma_f^2
\exp\!\left(
-\frac{1}{2}\left\|\frac{x_i-x_j}{\ell}\right\|_2^2
\right)
+
(\sigma_n^2+10^{-6})\delta_{ij}.
\]
The parameters are initialized as
$\ell=0.9$, $\sigma_f=1.1$, and $\sigma_n=0.15$. We use
$n=15{,}000$, and compare against a dense eigendecomposition reference for probe
value and gradient.

Figure~\ref{fig:trlogdet-adjoint} compares our gradient approximation with the
adjoint-based Lanczos implemented in \texttt{matfree}~\cite{kramer2024gradients} (referred to as Adjoint). We evaluate both methods
using the same forward Lanczos approximation, with and without full reorthogonalization.
With full reorthogonalization, our gradient approximation agrees closely with the
adjoint-based gradient once the forward Lanczos estimate is accurate. At \(m=60\), the
forward relative error is already below \(10^{-6}\), and the relative difference between
our gradient and the adjoint gradient is about \(3.4\times 10^{-5}\); at \(m=80\), these
drop to about \(10^{-11}\) and \(2.4\times 10^{-8}\), respectively. This supports the
central claim that, when the Lanczos estimate is accurate, our approximate gradients are accurate as well. We observe instability in the adjoint-based method after $m=100$. 
Without reorthogonalization, the
forward approximation converges more slowly and the adjoint gradient can become unstable,
whereas our approximation continues to track the forward accuracy and inherits the backward stability of the Lanczos quadrature.

This experiment focuses on a single probe
$u^\top\log(K(\theta))u$ with fixed $u$ sample. Additional probes could be used to improve the Hutchinson--Lanczos
log-determinant estimation. A detailed comparison for the full stochastic trace
estimator with GPyTorch's Lanczos-based log-determinant gradients, is provided in
Appendix~\ref{app:Gpy_logdet}.


\begin{figure}[t]
    \centering
    \includegraphics[width=0.6\linewidth]{figures/repo_forward_ours_vs_adjoint_n15000_d2_p1_seed0_refdense_gradient_comparison_notitle.png}
    \caption{
    Comparison of our method with the adjoint-based method~\cite{kramer2024gradients} for $u^\top \log K(\theta) u$ matrix experiment with $n=15{,}000$.
    Left: relative error of the Lanczos approximation to
    $u^\top\log(K(\theta))u$. Right: relative gradient error with respect to
    kernel hyperparameters.
    }
    \label{fig:trlogdet-adjoint}
\end{figure}





\subsection{Network sensitivity benchmarks}

We next evaluate our Lanczos gradient approximation on graph-sensitivity tasks based on
matrix exponentials. In applications, these sensitivities can
be used to rank influential edges or interactions, identify which parts of a network most
affect information flow, and assess the robustness of network statistics to structural
changes \citep{benzi2020,delacruz2022,pozza2018,schweitzer2023}. Following the
network-analysis setup of \citet{kressner2026}, we focus on sensitivities derived from
the matrix exponential on sparse graph adjacency matrices.

We consider two objectives. The first is the \emph{total-network-communicability} (TN)
sensitivity
\[
S^{(\mathrm{TN})}_{ij}(A)
=
e_i^\top L_{\exp}(A^\top,\mathbf{1}\mathbf{1}^\top)e_j,
\]
which measures how a global communicability score changes under perturbations of the
graph. The second is the \emph{subgraph-centrality} (SC) sensitivity
\[
S^{(\mathrm{SC})}_{ij}(A,\ell)
=
e_i^\top L_{\exp}(A^\top,e_\ell e_\ell^\top)e_j,
\]
which measures the sensitivity of the local subgraph-centrality score associated with
node $\ell$. The TN objective probes a dense global direction, whereas SC probes a
localized node direction. Together, they provide complementary tests of Fr\'echet-derivative
approximation, ranging from global network sensitivity to local node-level sensitivity.

In this setting, the Lanczos approximate gradient can be obtained very efficiently. Using~\eqref{eq:method_grad_approx}, and replacing the perturbation to be rank-$1$, $dA=ab^\top$, we get
\begin{equation}
    \operatorname{tr}\!\left(G V^\top a b^\top V\right)
    =
    (V^\top b)^\top G (V^\top a).
    \label{eq:exp_rank_one_contraction}
\end{equation}
Thus, the cost of computing the gradient is
$O(nm+m^2)$ without needing any backward pass. Both graph objectives above have this structure: TN uses the dense
rank-one direction $\mathbf{1}\mathbf{1}^\top$, while SC uses the sparse rank-one direction
$e_\ell e_\ell^\top$.

We compare our gradient approximation against the modified-Arnoldi
method of \citet{kressner2026} (referred to as KO), using a sparse block-exponential computation as the
reference for error measurement. We use their 
We run these experiments on sparse undirected graphs from the SNAP collection: \texttt{ca-HepTh}, an
arXiv high-energy-physics theory coauthorship network, and \texttt{com-Amazon}, an
Amazon product co-purchasing network with ground-truth communities. Additional dataset
details and exact details of our experiments for reproducibility are given in Appendix~\ref{app:network-sensitivity}.


Table~\ref{tab:network-sensitivity-summary} summarizes the graph-sensitivity
benchmarks at \(m=16\) and \(m=32\). Our method is consistently faster than KO across
both graphs and objectives, with average speedups from \(1.42\times\) on
\texttt{SNAP/ca-HepTh} to \(2.64\times\) on \texttt{com-Amazon}. This is consistent
with the Krylov problems being solved: our method works with the original \(n\times n\)
adjacency matrix, whereas KO uses a \(2n\times 2n\) block triangular embedding for the
Fr\'echet derivative action. By \(m=32\), both methods reach near machine-precision
relative gradient error; at \(m=16\), our method is more accurate except on
\texttt{ca-HepTh}/SC. Full curves are given in
Appendix~\ref{app:network-sensitivity}.



\begin{table}[t]
\centering
\small
\caption{
Graph-sensitivity benchmark summary at Krylov depths $m=16$ and $m=32$. Obj. defines the objective function tested, KO is the modified Arnoldi method in~\cite{kressner2026}.
Left: Average ($3$ iterations) runtime speedups over KO. Right: relative gradient errors with respect to the full Fr\'echet derivative.
}
\label{tab:network-sensitivity-summary}

\begin{minipage}[t]{0.36\linewidth}
\centering
\caption*{\small Speedup}
\begin{tabular}{llcc}
\toprule
Graph & Obj. & 16 & 32 \\
\midrule
HepTh & TN & $1.42$ & $1.55$ \\
Amazon & TN & $2.18$ & $2.58$ \\
HepTh & SC & $1.46$ & $1.60$ \\
Amazon & SC & $2.12$ & $2.64$ \\
\bottomrule
\end{tabular}
\end{minipage}
\hfill
\begin{minipage}[t]{0.60\linewidth}
\centering
\caption*{\small Relative gradient error}
\begin{tabular}{llcccc}
\toprule
& & \multicolumn{2}{c}{16} & \multicolumn{2}{c}{32} \\
\cmidrule(lr){3-4}\cmidrule(lr){5-6}
Graph & Obj. & Ours & KO & Ours & KO \\
\midrule
HepTh & TN & $10^{-4}$ & $10^{-4}$ & $10^{-14}$ & $10^{-11}$ \\
Amazon & TN & $10^{-7}$ & $10^{-5}$ & $10^{-14}$ & $10^{-12}$ \\
HepTh & SC & $10^{-1}$ & $10^{-2}$ & $10^{-10}$ & $10^{-12}$ \\
Amazon & SC & $10^{-8}$ & $10^{-8}$ & $10^{-14}$ & $10^{-15}$ \\
\bottomrule
\end{tabular}
\end{minipage}
\end{table}





\subsection{Quantum Hamiltonian learning}
\label{subsec:experiments_quantum}



Hamiltonian learning is a basic problem in quantum system identification: given
observations of a system's dynamics, the goal is to recover the Hamiltonian that
generates them. This task arises in the characterization and calibration of quantum
devices, validation of quantum simulators, and learning effective models for many-body
quantum systems \citep{banchi2023learning,huang2023robust,haah2024practical}. 
We consider a transverse-field Ising-type Hamiltonian on an \(L\)-site spin system. The
unknown Hamiltonian is represented as
\[
    H_\theta = \sum_{j=1}^p \theta_j P_j ,
\]
where \(P_j\in\mathbb{C}^{2^L\times 2^L}\) are local tensor-product operators. Let
\(I\), \(X\), and \(Z\) denote the \(2\times 2\) identity, Pauli-\(X\), and Pauli-\(Z\)
matrices. For \(i=0,\ldots,L-1\), define
\[
    X_i
    =
    I^{\otimes i}\otimes X\otimes I^{\otimes(L-i-1)},
    \qquad
    Z_i
    =
    I^{\otimes i}\otimes Z\otimes I^{\otimes(L-i-1)} .
\]
For \(i=0,\ldots,L-2\), define the nearest-neighbor interaction
\[
    Z_iZ_{i+1}
    =
    I^{\otimes i}\otimes Z\otimes Z\otimes I^{\otimes(L-i-2)} .
\]
The operator dictionary is
\[
    \{P_j\}_{j=1}^p
    =
    \{X_i\}_{i=0}^{L-1}
    \cup
    \{Z_i\}_{i=0}^{L-1}
    \cup
    \{Z_iZ_{i+1}\}_{i=0}^{L-2},
    \qquad
    p=3L-1 .
\]
In the reported experiment, \(L=8\), so the
Hilbert-space dimension is \(2^8=256\) and the number of unknown parameters is
\(p=23\).

We generate synthetic training data from a ground-truth Hamiltonian
\(H_{\theta_\star}\). The single-site \(X_i\) and \(Z_i\) coefficients are centered near
\(0.35\), while the nearest-neighbor \(Z_iZ_{i+1}\) coefficients are centered near
\(1.0\). Independent Gaussian perturbations with standard deviation \(0.05\) are added
to the coefficients. For each training example, we sample a random normalized initial
state \(x_j\in\mathbb{C}^{2^L}\), sample an evolution time
\[
    t_j \sim \mathrm{Unif}(2,6),
\]
and form the target state
\[
    c_j = \exp(-\mathrm{i}t_jH_{\theta_\star})x_j .
\]
The dataset contains \(30\) such input-output pairs. Given these samples, we recover the Hamiltonian parameters by minimizing
\[
    \mathcal{L}(\theta)
    =
    \sum_{j=1}^{30}
    \left\|
        c_j - \exp(-\mathrm{i}t_jH_\theta)x_j
    \right\|_2^2 .
\]
The optimization is initialized from a perturbation of the true parameter vector,
\[
    \theta_0 = \theta_\star + 0.1\,\xi,
    \qquad
    \xi \sim \mathcal{N}(0,I),
\]
and run for \(400\) steps with learning rate \(0.01\).

\begin{table}[t]
\centering
\small
\caption{
Hamiltonian-learning summary for $L=8$ with $400$ optimization steps.
Speedup is relative to the dense Fr\'echet-derivative baseline. The initial
gradient error is measured against the gradient with full eigendecomposition at $\theta_0$.
}
\label{tab:hamiltonian-learning}
\begin{tabular}{cccc}
\toprule
Method & Grad. rel. err. & Final param. error & Speedup \\
\midrule
Dense & -- & $2.00{\times}10^{-10}$ & $1.0{\times}$ \\
$m=30$ & $7.49{\times}10^{-2}$ & $5.43{\times}10^{-3}$ & $14.8{\times}$ \\
$m=40$ & $6.31{\times}10^{-4}$ & $1.61{\times}10^{-4}$ & $12.4{\times}$ \\
 $m=50$ & $5.62{\times}10^{-7}$ & $1.05{\times}10^{-7}$ & $10.4{\times}$ \\
\bottomrule
\end{tabular}
\end{table}


Table~\ref{tab:hamiltonian-learning} summarizes the Hamiltonian-learning experiment.
Increasing the Krylov depth sharply improves the gradient approximation at the
initial point: the relative gradient error drops from $7.49{\times}10^{-2}$ at
$m=30$ to $5.62{\times}10^{-7}$ at $m=50$. This translates directly into the
optimization behavior. With $m=30$, the method makes progress but plateaus away from
the dense solution; with $m=40$, it recovers the parameters to about $10^{-4}$; and
with $m=50$, it reaches $1.05{\times}10^{-7}$ parameter error while remaining
$10.4\times$ faster than the dense Fr\'echet baseline. Full loss and parameter-error
trajectories as functions of both iteration and wall time are shown in
Appendix~\ref{app:hamiltonian-learning}.







\section{Limitations and future work}
This work focuses on Hermitian matrices, where Lanczos yields a short three-term recurrence and a tridiagonal projected problem. Extending the approximation to general non-Hermitian matrices, for example via Arnoldi or structure-preserving Lanczos-type orthogonalization, is a natural direction for future work. Our analysis is also limited to scalar objectives \(u^\top f(A(\theta))v\). For derivatives of matrix-function-vector products \(f(A(\theta))v\), derivative-Krylov or adjoint-based methods may still be preferable. Our current formulation requires the algorithm to store the Lanczos vectors which can lead to a memory bottleneck for very large matrices and number of Lanczos iterations. Future work may reformulate the algorithm to reduce storage requirements.

Finally, our current implementation uses standard Python numerical and automatic-differentiation libraries, which introduce computational overhead. A compiled or accelerator-native implementation with careful batching and memory management is an important next step.

\section*{Acknowledgement}
This material is based upon work supported by the U.S. Department of Energy, Office of Science, Office of
Advanced Scientific Computing Research and Office of Basic Energy Sciences, Scientific Discovery through
Advanced Computing (SciDAC) program.

\bibliographystyle{plainnat}
\bibliography{sections/refs}

\appendix

\section{Proofs for Section~\ref{sec:grad_approx}}
\label{app:section4_details}

This appendix provides more details on the projected sensitivity,  and the residual-controlled error formula used in
Section~\ref{sec:grad_approx}. We use the same notation as in the main text.
The matrix \(A=A(\theta)\in\mathbb{R}^{n\times n}\) is symmetric and differentiable
with respect to the scalar parameter \(\theta\). The vector \(u\in\mathbb{R}^n\) is fixed
and nonzero. The Lanczos iteration is initialized at
\[
    v_1 = \frac{u}{\|u\|},
\]
and after \(m\) steps returns an orthonormal basis
\[
    V=[v_1,\ldots,v_m]\in\mathbb{R}^{n\times m}
\]
and a symmetric tridiagonal projected matrix
\[
    T = V^\top A V .
\]
We also use the Lanczos relation
\begin{equation}
    AV = VT+\beta_m v_{m+1}e_m^\top ,
    \label{eq:app_lanczos_relation}
\end{equation}
where \(e_m\) denotes the \(m\)-th canonical basis vector in \(\mathbb{R}^m\), i.e.,
the \(m\)-th column of the \(m\times m\) identity matrix. If \(\beta_m=0\), the Krylov subspace is
\(A\)-invariant and the residual term is zero.

The Lanczos estimate considered in the main text is
\begin{equation}
    \widehat{\phi}(\theta;u)
    =
    \|u\|^2 e_1^\top f(T)e_1 .
    \label{eq:app_lanczos_quad}
\end{equation}
All differentials below are with respect to \(\theta\), unless we explicitly view
\(\widehat{\phi}\) as a function of \(T\).

\subsection{Sensitivity with respect to the projected matrix}
\label{app:sensitivity}

Let
\[
    T = Q\Lambda Q^\top,
    \qquad
    \Lambda=\operatorname{diag}(\lambda_1,\ldots,\lambda_m),
\]
and define
\[
    c = Q^\top e_1 .
\]
Applying the symmetric Fréchet derivative formula
\eqref{eq:bg_hadamard_frechet} to the tridiagonal matrix \(T\), with perturbation \(dT\),
gives
\begin{equation}
    L_f(T,dT)
    =
    Q\bigl(F(\Lambda)\circ(Q^\top dTQ)\bigr)Q^\top ,
    \label{eq:app_frechet_formula}
\end{equation}
where \(F(\Lambda)\) is the divided-difference matrix defined in
\eqref{eq:bg_divided_difference}.

We now rewrite this scalar differential as a trace pairing with \(dT\). Since
\(c^\top M c=\operatorname{tr}(cc^\top M)\),
\begin{align}
    c^\top
    \bigl(F\circ(Q^\top dTQ)\bigr)c
    &=
    \operatorname{tr}\!\left(
        cc^\top \bigl(F\circ(Q^\top dTQ)\bigr)
    \right) \\
    &=
    \operatorname{tr}\!\left(
        \bigl((cc^\top)\circ F\bigr)^\top Q^\top dTQ
    \right).
\end{align}
The second equality uses the identity
\[
    \operatorname{tr}\!\left(A^\top(B\circ C)\right)
    =
    \operatorname{tr}\!\left((A\circ B)^\top C\right).
\]
Using cyclic invariance of the trace gives
\begin{align}
    d\widehat{\phi}
    &=
    \operatorname{tr}\!\left(
        \left[
        \|u\|^2 Q\bigl((cc^\top)\circ F\bigr)Q^\top
        \right]^\top dT
    \right).
\end{align}
Therefore
\begin{equation}
    d\widehat{\phi}
    =
    \operatorname{tr}(G^\top dT),
    \label{eq:app_dphi_G}
\end{equation}
with
\begin{equation}
    G
    =
    \|u\|^2 Q\bigl((cc^\top)\circ F\bigr)Q^\top .
    \label{eq:app_G}
\end{equation}
Since \(F\) is symmetric and \(cc^\top\) is symmetric, \(G\) is symmetric. This proves
Theorem~\ref{thm:projected_sensitivity}.


\subsection{Error from ignoring basis variation}
\label{app:error}

We now prove Theorem~\ref{thm:method_error}. The goal is to identify the exact
difference between the full differential and the forward-only approximation.

Let \(V^\perp\in\mathbb{R}^{n\times(n-m)}\) be an orthonormal basis for the orthogonal
complement of \(\operatorname{span}(V)\), chosen so that its first column is
\(v_{m+1}\). Since \(V^\top V=I\), differentiating the orthonormality constraint gives
\[
    (dV)^\top V+V^\top dV=0.
\]
Hence
\[
    S:=V^\top dV
\]
is skew-symmetric. Every \(dV\) can therefore be decomposed as
\begin{equation}
    dV = VS+V^\perp N,
    \label{eq:app_dV_decomp}
\end{equation}
where
\[
    S^\top=-S,
    \qquad
    N\in\mathbb{R}^{(n-m)\times m}.
\]
Define
\begin{equation}
    \eta = N^\top e_1^{(n-m)}\in\mathbb{R}^m,
    \label{eq:app_eta}
\end{equation}
where \(e_1^{(n-m)}\) is the coordinate vector corresponding to the \(v_{m+1}\)
direction in \(V^\perp\).

\begin{lemma}[Differential of the tridiagonal matrix]
\label{lem:app_dT}
Under the decomposition \eqref{eq:app_dV_decomp},
\begin{equation}
    dT
    =
    V^\top dA\,V
    +
    [T,S]
    +
    \beta_m(\eta e_m^\top+e_m\eta^\top),
    \label{eq:app_dT_decomp}
\end{equation}
where
\[
    [T,S]=TS-ST .
\]
\end{lemma}

\begin{proof}
Using the Lanczos relation \eqref{eq:app_lanczos_relation},
\[
    AV = VT+\beta_m v_{m+1}e_m^\top .
\]
Substituting \(dV=VS+V^\perp N\) into the first basis-variation term gives
\begin{align}
    dV^\top AV
    &=
    (VS+V^\perp N)^\top
    (VT+\beta_m v_{m+1}e_m^\top) \\
    &=
    S^\top T
    +
    \beta_m N^\top (V^\perp)^\top v_{m+1}e_m^\top \\
    &=
    S^\top T+\beta_m\eta e_m^\top .
\end{align}
Since \(S^\top=-S\),
\[
    dV^\top AV
    =
    -ST+\beta_m\eta e_m^\top .
\]
Similarly, transposing the Lanczos relation gives
\[
    V^\top A
    =
    TV^\top+\beta_m e_m v_{m+1}^\top .
\]
Hence
\begin{align}
    V^\top A\,dV
    &=
    \bigl(TV^\top+\beta_m e_m v_{m+1}^\top\bigr)(VS+V^\perp N) \\
    &=
    TS+\beta_m e_m\eta^\top .
\end{align}
Combining these two identities with
\[
    dT=dV^\top AV+V^\top dA\,V+V^\top A\,dV
\]
gives
\[
    dT
    =
    V^\top dA\,V
    +
    (TS-ST)
    +
    \beta_m(\eta e_m^\top+e_m\eta^\top),
\]
as claimed.
\end{proof}

Substituting Lemma~\ref{lem:app_dT} into the projected sensitivity formula
\eqref{eq:app_dphi_G} gives
\begin{align}
    d\widehat{\phi}
    &=
    \operatorname{tr}(G^\top V^\top dA\,V)
    +
    \operatorname{tr}(G^\top[T,S]) \\
    &\qquad
    +
    \beta_m
    \operatorname{tr}\!\left(
        G^\top(\eta e_m^\top+e_m\eta^\top)
    \right).
    \label{eq:app_error_decomp}
\end{align}
The first term is the forward-only direct term. It remains to show that the commutator
term vanishes under the fixed-starting-vector constraint.

Since \(u\) is fixed, the initial Lanczos vector \(v_1=u/\|u\|\) is fixed. Therefore
\[
    dv_1=0.
\]
The first column of \(dV=VS+V^\perp N\) gives
\[
    0=dv_1=VSe_1+V^\perp Ne_1.
\]
The two summands lie in orthogonal subspaces, so both must vanish:
\begin{equation}
    Se_1=0,
    \qquad
    Ne_1=0.
    \label{eq:app_fixed_start}
\end{equation}

\begin{lemma}[Vanishing commutator]
\label{lem:app_commutator}
If \(S^\top=-S\) and \(Se_1=0\), then
\[
    \operatorname{tr}(G^\top[T,S])=0 .
\]
\end{lemma}

\begin{proof}
Let
\[
    \widetilde S = Q^\top S Q .
\]
Since \(S^\top=-S\), we also have \(\widetilde S^\top=-\widetilde S\). Moreover,
using \(c=Q^\top e_1\) and \(Se_1=0\),
\[
    \widetilde S c
    =
    Q^\top S Q Q^\top e_1
    =
    Q^\top S e_1
    =
    0 .
\]
Because \(T=Q\Lambda Q^\top\),
\[
    Q^\top[T,S]Q
    =
    Q^\top(TS-ST)Q
    =
    \Lambda \widetilde S-\widetilde S\Lambda .
\]
Using the definition of \(G\) from \eqref{eq:app_G} and cyclic invariance of the trace,
\begin{align}
    \operatorname{tr}(G^\top[T,S])
    &=
    \|u\|^2
    \operatorname{tr}\!\left(
        \bigl((cc^\top)\circ F(\Lambda)\bigr)
        (\Lambda\widetilde S-\widetilde S\Lambda)
    \right).
\end{align}
It remains to show that the trace on the right-hand side is zero. Define
\[
    \mathcal{C}
    =
    \operatorname{tr}\!\left(
        \bigl((cc^\top)\circ F(\Lambda)\bigr)
        (\Lambda\widetilde S-\widetilde S\Lambda)
    \right).
\]
Expanding entrywise,
\begin{align}
    \mathcal{C}
    &=
    \sum_{i,j}
    c_i c_j F(\lambda_i,\lambda_j)
    (\lambda_j-\lambda_i)\widetilde S_{ji}.
    \label{eq:app_commutator_expansion}
\end{align}
The diagonal terms vanish because \(\lambda_j-\lambda_i=0\) when \(i=j\). For
\(i\neq j\), the divided-difference identity gives
\[
    F(\lambda_i,\lambda_j)(\lambda_j-\lambda_i)
    =
    f(\lambda_j)-f(\lambda_i).
\]
Hence
\begin{align}
    \mathcal{C}
    &=
    \sum_{i,j}
    c_i c_j
    \bigl(f(\lambda_j)-f(\lambda_i)\bigr)
    \widetilde S_{ji} \\
    &=
    \sum_{i,j}
    c_i c_j f(\lambda_j)\widetilde S_{ji}
    -
    \sum_{i,j}
    c_i c_j f(\lambda_i)\widetilde S_{ji}.
\end{align}
The two sums can be written in matrix form as
\[
    \mathcal{C}
    =
    (f(\Lambda)c)^\top \widetilde S c
    -
    c^\top \widetilde S f(\Lambda)c .
\]
Since \(\widetilde S\) is skew-symmetric,
\[
    c^\top \widetilde S f(\Lambda)c
    =
    - (f(\Lambda)c)^\top \widetilde S c .
\]
Therefore
\[
    \mathcal{C}
    =
    2(f(\Lambda)c)^\top \widetilde S c .
\]
Finally, \(\widetilde S c=0\), so \(\mathcal{C}=0\). Thus
\[
    \operatorname{tr}(G^\top[T,S])=0 .
\]
\end{proof}
We now complete the proof of Theorem~\ref{thm:method_error}. By
Lemma~\ref{lem:app_commutator}, the commutator term in
\eqref{eq:app_error_decomp} is zero. Hence
\[
    d\widehat{\phi}
    =
    \operatorname{tr}(G^\top V^\top dA\,V)
    +
    \beta_m
    \operatorname{tr}\!\left(
        G^\top(\eta e_m^\top+e_m\eta^\top)
    \right).
\]
The direct term can be rewritten as
\[
    \operatorname{tr}(G^\top V^\top dA\,V)
    =
    \operatorname{tr}\!\left((VGV^\top)^\top dA\right).
\]
Since \(G\) is symmetric, this is also
\[
    \operatorname{tr}\!\left(VGV^\top dA\right)
\]
for symmetric \(dA\), as in \eqref{eq:method_error_formula}. Finally,
\begin{align}
    \operatorname{tr}\!\left(
        G^\top(\eta e_m^\top+e_m\eta^\top)
    \right)
    &=
    \operatorname{tr}(G\eta e_m^\top)
    +
    \operatorname{tr}(G e_m\eta^\top) \\
    &=
    e_m^\top G\eta+\eta^\top G e_m \\
    &=
    2e_m^\top G\eta .
\end{align}
Therefore
\begin{equation}
    d\widehat{\phi}
    =
    \operatorname{tr}\!\left(VGV^\top dA\right)
    +
    2\beta_m e_m^\top G\eta ,
    \label{eq:app_method_error_formula}
\end{equation}
which is precisely \eqref{eq:method_error_formula}. 

\section{Additional experimental results}
\label{app:Extra_exp_results}
\subsection{Log-determinant optimization in GP training}
\label{app:Gpy_logdet}
We evaluate the proposed gradient approximation in an end-to-end exact
Gaussian-process training problem. Given training data
\(\{(x_i,y_i)\}_{i=1}^{n}\), we optimize the negative log marginal likelihood
\[
    \mathcal{L}(\theta)
    =
    \frac{1}{2n}(y-\mu_\theta)^\top K_\theta^{-1}(y-\mu_\theta)
    +
    \frac{1}{2n}\log\det K_\theta
    +
    \frac{1}{2}\log(2\pi),
\]
where
\[
    K_\theta
    =
    \sigma_f^2 K_{\nu=3/2}(X,X;\ell)
    +
    \sigma_n^2 I.
\]
Here \(K_{\nu=3/2}\) is a Mat\'ern-\(3/2\) kernel with ARD lengthscales,
\(\sigma_f^2\) is the output scale, and \(\sigma_n^2\) is the observation noise.

The inverse-quadratic term is computed using GPyTorch's conjugate-gradient
solver in both methods. The difference is the treatment of the log-determinant
gradient. GPyTorch uses its standard stochastic Lanczos quadrature estimator and
its built-in backward approximation. Our method instead replaces the
log-determinant gradient with the proposed projected Lanczos-gradient estimator.

We use a rank-\(r\) pivoted-Cholesky preconditioner
\[
    P_\theta = R_\theta R_\theta^\top + \delta_\theta I
\]
and apply the identity
\[
    \log\det K_\theta
    =
    \log\det P_\theta
    +
    \log\det\!\left(P_\theta^{-1/2}K_\theta P_\theta^{-1/2}\right).
\]
For each probe vector \(z_s\), we run \(m\) steps of unreorthogonalized Lanczos
on
\[
    B_\theta = P_\theta^{-1/2}K_\theta P_\theta^{-1/2},
\]
starting from \(z_s/\|z_s\|\), producing a tridiagonal matrix \(T_s\). The
log-determinant is estimated as
\[
    \widehat{\log\det K_\theta}
    =
    \log\det P_\theta
    +
    \frac{n}{p}\sum_{s=1}^{p}
    e_1^\top \log(T_s)e_1 .
\]
For the gradient, we use our gradient approximation, and evaluate the resulting matrix-vector contractions using automatic
differentiation through \(K_\theta\) and \(P_\theta\). Thus the derivative of the
preconditioner with respect to the GP hyperparameters is included via autodifferentiation. The preconditioner is included here as it is crucial for improving the conditioning of the matrix thus improving convergence of CG and Lanczos.

We use the Protein dataset from the UCI benchmark collection~\citep{rana2013protein}.
We select a seeded
subset of \(4096\) examples, with an \(80/20\) train/test split, giving
\(3276\) training points and \(820\) test points. Both methods use the same
initial hyperparameters, Adam optimizer with learning rate \(0.025\), double
precision, \(p=20\) probes, \(m=40\) Lanczos steps, and a rank-15
pivoted-Cholesky preconditioner. Test RMSE is evaluated every three iterations.

Figure~\ref{fig:gp-protein-training} shows the training trajectory. GPyTorch
decreases the stochastic training objective more rapidly, but its test RMSE
bottoms out early and then increases which shows overfitting of parameters. Our Lanczos-gradient
method decreases the training objective more slowly, but gives a steadily
improving test RMSE and reaches the best final predictive accuracy in this run. 

To isolate the log-determinant approximation at initialization, we compare the
first-iteration stochastic log-determinant estimate and its gradient against a
dense Cholesky reference. For this diagnostic, both methods use the same initial
hyperparameters, the same rank-15 pivoted-Cholesky preconditioner, and the same
set of Gaussian probe vectors. Table~\ref{tab:gp-first-iter-logdet} shows that
the two methods obtain comparable log-determinant estimates, while our
preconditioned non-reorthogonalized Lanczos estimator gives a smaller gradient
error in this run. Computing dense references throughout training, while also ensuring that all
stochastic estimators use identical probe vectors at every iteration, is
expensive and application-specific. We therefore leave a systematic study of how
log-determinant gradient accuracy affects downstream GP training to future work.



While a more systematic study of optimization stability, preconditioning choices,
and hyperparameter trajectories is left for future work, this experiment
demonstrates that the proposed gradient estimator can be competitive in
end-to-end practical Gaussian-process training.

\begin{table}[t]
\centering
\small
\caption{First-iteration log-determinant diagnostic for exact-GP training on
Protein. Both stochastic methods use the same initial hyperparameters, same
Gaussian probes, \(m=40\), \(p=20\), and a rank-15 pivoted-Cholesky
preconditioner. Errors are relative to a dense Cholesky reference.}
\label{tab:gp-first-iter-logdet}
\begin{tabular}{lcc}
\toprule
Method & Rel. logdet error & Rel. gradient error \\
\midrule
GPyTorch & $2.29\times 10^{-3}$ & $1.44\times 10^{-1}$ \\
Ours & $1.94\times 10^{-3}$ & $5.10\times 10^{-3}$ \\
\bottomrule
\end{tabular}
\end{table}




\begin{figure}[t]
    \centering
    \includegraphics[width=\linewidth]{figures/GP_loss_and_RMSE.png}
    \caption{GP training on the Protein dataset with a Mat\'ern-\(3/2\)
    kernel. Both methods use \(n=4096\), \(m=40\), \(p=20\), and a rank-15
    pivoted-Cholesky preconditioner. GPyTorch decreases the stochastic marginal
    likelihood faster, while our unreorthogonalized Lanczos-gradient method
    reaches lower test RMSE.}
    \label{fig:gp-protein-training}
\end{figure}




\subsection{Graph-sensitivity results}
\label{app:network-sensitivity}

We evaluate the proposed forward-only gradient approximation on graph-sensitivity
objectives involving the matrix exponential. Let \(A\) be the adjacency matrix of an
undirected graph. We consider two standard matrix-function centrality objectives.

The first is total network communicability,
\[
    \mathrm{TN}(A)
    =
    \mathbf{1}^\top \exp(A)\mathbf{1}.
\]
The sensitivity of \(\mathrm{TN}\) with respect to an entrywise perturbation of \(A\) is
\[
    S^{\mathrm{TN}}_{ij}(A)
    =
    \mathbf{1}^\top L_{\exp}(A,e_i e_j^\top)\mathbf{1}.
\]
Using the adjoint identity for Fr\'echet derivatives, this can equivalently be written as
\[
    S^{\mathrm{TN}}_{ij}(A)
    =
    e_i^\top L_{\exp}(A^\top,\mathbf{1}\mathbf{1}^\top)e_j .
\]

The second is subgraph centrality at node \(\ell\),
\[
    \mathrm{SC}_{\ell}(A)
    =
    e_\ell^\top \exp(A)e_\ell .
\]
Its entrywise sensitivity is
\[
    S^{\mathrm{SC}}_{ij}(A;\ell)
    =
    e_\ell^\top L_{\exp}(A,e_i e_j^\top)e_\ell,
\]
or equivalently
\[
    S^{\mathrm{SC}}_{ij}(A;\ell)
    =
    e_i^\top L_{\exp}(A^\top,e_\ell e_\ell^\top)e_j .
\]
Thus both TN and SC sensitivities reduce to Fr\'echet derivative actions with a
structured direction matrix: \(\mathbf{1}\mathbf{1}^\top\) for TN and
\(e_\ell e_\ell^\top\) for SC.

We compare our Lanczos-based forward-only gradient estimator against the modified
Arnoldi method of \citet{kressner2026}. For \texttt{ca-HepTh}, we use
\((i,j)=(7200,6969)\), following the network-sensitivity setup of
\citet{kressner2026}. For \texttt{com-Amazon}, we use \((i,j)=(0,53525)\).
For SC, we set \(\ell=7200\) on \texttt{ca-HepTh} and \(\ell=0\) on
\texttt{com-Amazon}. In all experiments, we use the raw undirected adjacency
matrix and \(f(A)=\exp(A)\).

Figure~\ref{fig:network-sensitivity-speedup} reports speedups at Krylov depths
\(m=16\) and \(m=32\), computed as the modified-Arnoldi runtime divided by our
runtime and averaged over three runs. Our method is consistently faster across
both graphs and both objectives. 

Figures~\ref{fig:network-sensitivity-grad-error} and
\ref{fig:network-sensitivity-value-error} report relative gradient and value
errors over the full Krylov-depth sweep. By \(m=32\), both methods reach very
small relative gradient error on all benchmarks. Our method gives lower gradient
error on the TN objectives and on \texttt{com-Amazon}/SC, while modified Arnoldi
is more accurate on \texttt{ca-HepTh}/SC at intermediate depths. For the function
value itself, our Lanczos approximation is consistently more accurate across the
tested depths. While the accuracy difference between the two is not much, our method is much more efficient in terms of computation as it does not use reorthogonalization and works with a smaller matrix.

\begin{figure}[t]
    \centering
    \includegraphics[width=0.62\linewidth]{figures/paper_batched_safe_r3_speedup_mean_ci.png}\
    \caption{Speedup of the proposed method over modified Arnoldi on graph-sensitivity benchmarks. Bars show mean speedup over three runs; error bars show 95\% confidence intervals.}
    \label{fig:network-sensitivity-speedup}
\end{figure}

\begin{figure}[t]
    \centering
    \includegraphics[width=0.7\linewidth]{figures/Combined_gradient_error.png}
    \caption{Relative gradient error versus Krylov depth \(m\) for TN and SC sensitivities on \texttt{ca-HepTh} and \texttt{com-Amazon}.}
    \label{fig:network-sensitivity-grad-error}
\end{figure}

\begin{figure}[t]
    \centering
    \includegraphics[width=0.7\linewidth]{figures/Combined_value_error.png}
    \caption{Relative value error versus Krylov depth \(m\) for TN and SC bilinear forms on \texttt{ca-HepTh} and \texttt{com-Amazon}.}
    \label{fig:network-sensitivity-value-error}
\end{figure}







\begin{figure}[t]
\centering
\begin{subfigure}[t]{0.48\linewidth}
  \centering
  \includegraphics[width=\linewidth]{figures/loss_vs_step.png}
  \caption{Training loss vs. optimization step.}
\end{subfigure}
\hfill
\begin{subfigure}[t]{0.48\linewidth}
  \centering
  \includegraphics[width=\linewidth]{figures/param_error_vs_step.png}
  \caption{Parameter error vs. optimization step.}
\end{subfigure}
\caption{
Hamiltonian-learning convergence as a function of optimization step. Larger Krylov
depths more closely track the dense Fr\'echet-derivative baseline.
}
\label{fig:hamiltonian-learning-step}
\end{figure}

\subsection{Hamiltonian-learning results}
\label{app:hamiltonian-learning}

Figure~\ref{fig:hamiltonian-learning-step} shows the training loss and parameter
recovery error as functions of optimization step. The Krylov
gradient approximation exhibits a clear accuracy--cost tradeoff as the Lanczos depth
is varied. For $m=30$, the optimization initially follows the dense trajectory but
then plateaus, consistent with a non-negligible gradient approximation error. For
$m=40$, the method tracks the dense baseline for much longer and reaches a substantially
smaller final error. For $m=50$, the Krylov run nearly overlaps the dense loss and
parameter-error curves as a function of iteration, while converging much faster in time.

These results complement Table~\ref{tab:hamiltonian-learning}: increasing $m$ improves
the quality of the gradient approximation and therefore the attainable optimization
accuracy, while the Krylov methods remain substantially faster than the dense
Fr\'echet-derivative computation.


% \input{checklist}


\end{document}
